\subsection{Ejercicio 5.3: análisis del desempeño de YYC y BT}
\label{sec:analisis53}

El objetivo de este apartado es comparar los tiempos de YYC y BT y observar
cómo cambia su desempeño al ordenar las filas por cantidad ascendente de
unos. Se utilizan las mediciones manuales registradas durante las ejecuciones;
Los tiempos de BT corresponden a la versión histórica identificada en 5.2,
no a la implementación corregida integrada posteriormente. Tampoco
se mezclan con la serie automática previa que tenía un límite de
60 segundos por proceso.

\subsubsection*{Condiciones de las mediciones}

La serie local se realizó en un equipo con procesador Intel Core Ultra 7
155H, 16 GB de RAM y Windows de 64 bits. La frecuencia de 1.40 GHz mostrada
en la información del sistema es un dato de identificación del equipo, no
una medición de la frecuencia efectiva durante cada ejecución. Estas
especificaciones corresponden a esta serie del ejercicio 5 y son distintas
del entorno descrito para los ejercicios 1 a 3.

Se registró una ejecución por combinación de matriz, algoritmo y orden de
filas. El orden ascendente se obtuvo contando los unos de cada fila; en
caso de empate se conservó su orden relativo. El ordenamiento se realizó
antes de ejecutar el algoritmo, por lo que su costo no se incluye en los
tiempos presentados. Todos los valores de las tablas se expresan en
milisegundos: los microsegundos se dividieron para 1000 y los segundos se
multiplicaron por 1000.

Se tomó el valor que el programa muestra como tiempo total de ejecución.
En YYC corresponde a la suma de los intervalos de procesamiento de las
filas; en BT corresponde al intervalo de resolución del algoritmo, sin la
verificación posterior por minimalidad. La impresión y exportación de
resultados quedan fuera de estos intervalos. Los criterios de medición
son, por tanto, los implementados en cada algoritmo y no el tiempo total
que permanece abierta la terminal.

También se recibió un registro de BT sobre $\gamma^3(\theta(A,B))$, en orden
original, de 19 horas, 14 minutos y 43 segundos, reportado por el programa:
\[
19\cdot3600+14\cdot60+43=69\,283\text{ s}
=69\,283\,000\text{ ms}.
\]
Este dato se identifica con un asterisco porque no se verificó la
equivalencia de sus condiciones de ejecución con las de la serie local.
Se conserva como registro adicional, sin atribuirlo al equipo local ni
utilizarlo para calcular ventajas de un algoritmo sobre otro.

Como se dispone de una observación por caso y no se controlaron de forma
documentada la carga del sistema, la temperatura o la posible concurrencia
de procesos, las diferencias describen estas ejecuciones. No permiten
estimar una variabilidad estadística ni concluir que diferencias pequeñas
se repetirán en cualquier sesión.

\subsubsection*{Resultados registrados}

Se utiliza $C=\theta(A,B)$. Las mediciones analizadas son las de
las Tablas~\ref{tab:resultadosphi52} y~\ref{tab:resultadosgamma52}
del punto 5.2. Los casos sin resultado final conservan el estado
\emph{falta ejecutar}; sus proyecciones se presentan separadas
más adelante y no se incorporan a las tablas de tiempos medidos.

\subsubsection*{Comparación entre los algoritmos}

En los casos de la serie local donde ambos algoritmos finalizaron, YYC
registró un tiempo menor que BT. La diferencia se hace especialmente
grande en $\varphi^3(C)$: en orden original, YYC tardó 231.593 ms,
mientras BT tardó 806\,341 ms, aproximadamente 13 minutos y 26 segundos.
La razón entre esos tiempos es cercana a 3482. En orden ascendente,
YYC tardó 57.589 ms y BT 3\,055\,669 ms, aproximadamente 50 minutos y
56 segundos; en esa ejecución, la razón fue cercana a 53\,060.
Estas razones se refieren a las mediciones registradas y no a constantes
de rendimiento de los algoritmos.

La diferencia tiene una explicación en sus estrategias de búsqueda.
YYC construye los testores típicos de manera incremental y conserva los
conjuntos compatibles de las filas ya procesadas. BT recorre candidatos
en orden binario y aplica saltos para omitir bloques que no pueden
contener testores típicos. Sin embargo, durante ese recorrido también
puede analizar testores que no son típicos. Por ello, tener pocos
testores típicos finales no garantiza que el recorrido de BT sea corto.

En la familia $\varphi^N(C)$ se mantienen 48 filas y las columnas pasan
de 14 a 70. La cantidad de testores típicos es
\[
5N^2+7N^3+2N^4,
\]
pero el espacio de todos los subconjuntos de columnas tiene
$2^{14N}$ elementos, incluido el vacío. Por ejemplo, con 70 columnas
existen $2^{70}\approx1.18\cdot10^{21}$ subconjuntos posibles. BT no
los visita necesariamente todos: sus saltos reducen la búsqueda. Aun
así, esta escala ayuda a comprender por qué algunas entradas pueden
requerir mucho trabajo aunque su matriz tenga pocas filas.

En la familia $\gamma^N(C)$ aumentan tanto las filas como las columnas:
la matriz tiene $48N$ filas, $14N$ columnas y $14^N$ testores típicos.
Esto incrementa también el volumen de resultados que cualquier
algoritmo completo debe producir. Las dos familias permiten estudiar
fenómenos distintos, como propone la metodología de \cite{alba2016}.
No se puede explicar el tiempo utilizando solamente la cantidad de
filas, columnas o testores finales de manera aislada.

\subsubsection*{Efecto del ordenamiento sobre YYC}

En las cuatro matrices $\varphi^2(C),\ldots,\varphi^5(C)$, el orden
ascendente produjo un tiempo menor que el original. Para expresar
esa diferencia se calcula el factor observado
\[
S=\frac{t_{\mathrm{original}}}{t_{\mathrm{ascendente}}}.
\]
Los factores fueron aproximadamente 2.63, 4.02, 5.03 y 6.43,
respectivamente. Por ejemplo, en $\varphi^5(C)$ el tiempo pasó de
14.095 s a 2.192 s.

Una explicación posible es que las filas con menos unos imponen
restricciones más fuertes al inicio y pueden reducir la cantidad de
conjuntos intermedios que YYC debe mantener o extender. Esta es una
interpretación consistente con su funcionamiento; no se midió en esta
serie un perfil completo de operaciones que permita atribuir toda
la diferencia a una sola causa.

El efecto no fue uniforme en la familia $\gamma$. En $\gamma^2(C)$
el tiempo bajó de 34.407 a 11.666 ms. En cambio, en $\gamma^3(C)$
subió de 280.705 a 337.262 ms, un aumento aproximado del 20.15\%,
y en $\gamma^4(C)$ pasó de 10.640 a 15.790 s, un aumento del 48.40\%.
Por tanto, ordenar por cantidad de unos no es una regla que siempre
acelere YYC. También influye cómo se distribuyen esas entradas entre
columnas y cómo las filas de distintos bloques interactúan durante
la construcción de los conjuntos intermedios.

\subsubsection*{Efecto del ordenamiento sobre BT}

Los resultados completos de la serie local tampoco muestran una
mejora uniforme para BT. En $\varphi^2(C)$ el tiempo pasó de 1.169
a 1.124 s, una diferencia pequeña en una sola observación.
En $\varphi^3(C)$ sucedió lo contrario: aumentó de 806.341 a
3055.669 s, aproximadamente 3.79 veces. En $\gamma^2(C)$ disminuyó
de 31.476 a 25.520 s. En el caso base $C$ aumentó de 0.7632 a
1.027 ms, aunque su corta duración aconseja no generalizar a partir
de esa diferencia aislada.

En la versión histórica de BT usada para estas mediciones, el salto
aplicado a un candidato que no es testor depende
de la primera fila que queda sin cubrir. Cambiar el orden de las
filas puede cambiar esa primera fila y, con ello, la posición que
determina el salto y el trabajo de comprobación. Ordenar por número
de unos no optimiza directamente ese mecanismo. Esto explica por
qué la misma ordenación puede ayudar a un algoritmo y perjudicar
al otro. En particular, en $\varphi^3(C)$ favoreció a YYC y coincidió
con un aumento considerable del tiempo de esa versión de BT.

En la implementación corregida, S1 considera todas las filas no cubiertas
y utiliza el último uno situado más a la izquierda. Reordenar las mismas
filas no cambia ese límite ni la secuencia de candidatos; sí puede variar
el costo de las comprobaciones. Los factores de tiempo de la versión
histórica no se trasladan a la implementación corregida.

\subsubsection*{Falta ejecutar}

Los siguientes casos no tienen un resultado final registrado en
la serie utilizada para este reporte:
\begin{itemize}[nosep]
\item BT original y ascendente sobre $\varphi^4(C)$.
\item BT original y ascendente sobre $\varphi^5(C)$.
\item BT ascendente sobre $\gamma^3(C)$.
\item BT original y ascendente sobre $\gamma^4(C)$.
\end{itemize}
Su estado es \textbf{falta ejecutar}. No se analizan como resultados
obtenidos ni se les asigna un tiempo observado.

\subsubsection*{Estimaciones exploratorias de los casos pendientes}

Para describir la escala que podrían alcanzar las ejecuciones se
plantea un \textbf{escenario de extrapolación geométrica no validado}.
El escenario corresponde exclusivamente a la versión histórica de BT.
Estas cifras no sustituyen los casos pendientes ni predicen de forma
fiable cuándo terminará un proceso que ya está ejecutándose. Tampoco
son cotas inferiores, superiores ni intervalos de confianza.

Para cada orden de filas de la familia $\varphi$, se toma el factor
entre los dos últimos casos completos de BT:
\[
r_{\varphi,o}=\frac{806.341}{1.169}\approx689.77,
\qquad
r_{\varphi,a}=\frac{3055.669}{1.124}\approx2718.57.
\]
Todos los valores de estas razones están en segundos. Si, como
hipótesis, se mantuviera ese factor por cada copia adicional,
el modelo sería
\[
\widehat t_{\varphi,o}(N)=806.341\,r_{\varphi,o}^{N-3},
\qquad
\widehat t_{\varphi,a}(N)=3055.669\,r_{\varphi,a}^{N-3}.
\]

Para $\gamma$ puede construirse únicamente un escenario
\emph{condicional}: suponiendo que el registro prolongado sea
comparable con el de $\gamma^2(C)$ original, el factor sería
\[
r_\gamma=\frac{69\,283}{31.476}\approx2201.14.
\]
Aplicar ese mismo factor al orden ascendente requiere una segunda
hipótesis, aún no comprobada. Bajo ambas hipótesis se tendría
$\widehat t_{\gamma,o}(4)=69\,283r_\gamma$,
$\widehat t_{\gamma,a}(3)=25.520r_\gamma$ y
$\widehat t_{\gamma,a}(4)=25.520r_\gamma^2$.

\begin{table}[htbp]
\centering
\small
\begin{tabular}{l l r}
\toprule
Matriz & Orden de BT & Proyección del escenario\\
\midrule
$\varphi^4(C)$ & Original   & 6.44 días\\
$\varphi^4(C)$ & Ascendente & 96.15 días\\
$\varphi^5(C)$ & Original   & 12.16 años\\
$\varphi^5(C)$ & Ascendente & 715.62 años\\
$\gamma^3(C)$  & Ascendente & 15.60 horas\\
$\gamma^4(C)$  & Original   & 4.83 años\\
$\gamma^4(C)$  & Ascendente & 3.92 años\\
\bottomrule
\end{tabular}
\caption{Estimaciones exploratorias, no tiempos medidos ni pronósticos
validados. Para expresar años se usaron 365.25 días. Las proyecciones
de $\gamma$ dependen además de una equivalencia de condiciones no verificada.}
\label{tab:escenario53}
\end{table}

Los valores extremos evidencian la sensibilidad del modelo a sus
supuestos. Por ejemplo, la cifra de 715.62 años no significa que
una ejecución vaya a durar ese tiempo: resulta de prolongar dos
veces una razón muy grande obtenida con sólo dos observaciones.
No hay datos suficientes para validar que esa razón permanezca
constante. La conclusión defendible es que algunos casos pueden
ser imprácticos de completar dentro del plazo disponible, no que
se conozca su fecha de finalización. Hasta obtener una ejecución
completa, conservan el estado \emph{falta ejecutar}.

La explicación de las duraciones elevadas no debe reducirse a
calificar el computador como lento. La implementación de BT
recorre sus candidatos de forma secuencial; disponer de varios
núcleos no convierte automáticamente ese recorrido en paralelo.
Además, un equipo con 16 GB de RAM no elimina el crecimiento
combinatorio del trabajo. La potencia del procesador y la carga
del sistema pueden modificar el tiempo, pero la estructura de
la matriz y la estrategia de búsqueda siguen siendo determinantes.

\subsubsection*{Conclusiones del análisis}

En las comparaciones completas de la serie local, YYC obtuvo
menores tiempos que BT. Ordenar las filas por cantidad ascendente
de unos benefició a YYC en la familia $\varphi$, pero no en todos
los casos de $\gamma$. En BT se observaron tanto disminuciones
como aumentos, incluido el incremento de aproximadamente 3.79 veces
en $\varphi^3(C)$. Por ello, el ordenamiento debe evaluarse para
cada familia de matrices y no asumirse como una mejora universal.

Los resultados apoyan la utilidad de experimentar con distintas
estructuras de matrices, en lugar de concluir el desempeño de un
algoritmo a partir de una única entrada. Las conclusiones anteriores
se limitan a las ejecuciones registradas; los casos pendientes y
las proyecciones no se utilizan para afirmar resultados que todavía
no se han obtenido.
